\chapter{Оператор числа и целочисленная кратность проходов одной нити}
\label{ch:v6-single-thread-c4-suspension-parent-gate}

\section{Коррекция переменной}

Предыдущий гейт сформулировал возможное спасение через разные времена
пребывания на четырёх тетраэдральных направлениях. Для исходной картины
одной нити это неточно. Фундаментальной величиной должна быть не
непрерывная длительность, а целое число полных проходов через голономный
цикл:
\begin{equation}
 K_a\in\mathbb N.
\end{equation}

Ключевая зацепка уже содержится в неограниченном родителе Тёплица Тома V
\cite{single-thread-bourne-kellendonk-rennie2017,
single-thread-arici-mesland2020}:
\begin{equation}
 Ne_n=ne_n,
 \qquad
 Ue_n=e_{n+1},
 \qquad
 [N,U]=U.
\end{equation}
Следовательно,
\begin{equation}
 [N,U^k]=kU^k.
\end{equation}
Оператор $N$ различает первый, миллиардный и десятимиллиардный витки без
перехода к непрерывному времени.

\section{Голономия хранит остаток, а не полное число}

Конечная голономия видит только класс числа проходов:
\begin{equation}
 k\longmapsto
 \bigl(k\bmod4,\;k\bmod3,\;k\bmod6\bigr)
\end{equation}
для проектных $C_4$, $Z_3$ и $Z_6$. Например,
\begin{equation}
 10^{10}\bmod4=0,
 \qquad
 10^{10}\bmod3=1,
 \qquad
 10^{10}\bmod6=4.
\end{equation}
Но при этом собственное значение $N$ остаётся равным $10^{10}$, а индекс
компрессии $PU^{10^{10}}P$ равен $-10^{10}$.

Поэтому роль $C_4$ уточняется принципиально:
\begin{quote}
$C_4$ не обязан переставлять четыре тетраэдральные оси. Он является
остаточной фазой абсолютного счётчика $N$ по модулю четыре.
\end{quote}
Тем самым нечётность четырёхцикла из предыдущего гейта остаётся верной для
буквальной перестановки осей, но не запрещает подъём на целочисленную
лестницу проходов.

\section{Один цикл с заданными кратностями}

Пусть через локальную область нить проходит по четырём направлениям
$K_0,K_1,K_2,K_3$ раз. Заменим четыре грубые метки на
\begin{equation}
 M=K_0+K_1+K_2+K_3
\end{equation}
различимых позиций обхода и соединим их правилом
\begin{equation}
 m\longmapsto m+1\pmod M.
\end{equation}

\begin{theorem}[Целочисленный одноконтурный подъём]
Для любого строго положительного вектора кратностей
$(K_0,K_1,K_2,K_3)$ существует один замкнутый цикл из $M$ позиций, в
котором каждая позиция имеет ровно одного предшественника и одного
последователя, а направление $a$ встречается ровно $K_a$ раз.
\end{theorem}

Это снимает прежний комбинаторный запрет $3+1$: он относился к четырём
меткам как к четырём состояниям перестановки. После подъёма каждая метка
имеет множество отдельных посещений, и все они могут лежать на одной
циклической нити.

Однако сами числа $K_a$ не выбирают порядок букв в циклическом слове.
Разные слова имеют одинаковые частоты, но могут задавать разные узлы,
соседства и локальные напряжения.

\section{Приближение рабочего вакуума целыми проходами}

Рабочие доли локальных проходов равны
\begin{equation}
 w=(0.9011874253513867,
 0.032937524882871085^{(3)}).
\end{equation}
Для симметричных целых кратностей $(K_0,K_q,K_q,K_q)$ получены
последовательные рациональные приближения:

\begin{center}
\begin{tabular}{r|r|r|r|r}
$K_0$ & $K_q$ & $M$ & $\max|K_a/M-w_a|$ & остаток второго момента\\
\hline
$219$ & $8$ & $243$ & $4.71\cdot10^{-5}$ & $5.13\cdot10^{-5}$\\
$2353$ & $86$ & $2611$ & $1.41\cdot10^{-7}$ & $1.53\cdot10^{-7}$\\
$15103$ & $552$ & $16759$ & $3.67\cdot10^{-9}$ & $3.99\cdot10^{-9}$\\
$88265$ & $3226$ & $97943$ & $1.22\cdot10^{-11}$ & $1.32\cdot10^{-11}$\\
$2243206$ & $81987$ & $2489167$ & $1.47\cdot10^{-13}$ & $1.59\cdot10^{-13}$\\
$22696855$ & $829548$ & $25185499$ & $3.00\cdot10^{-15}$ & $3.13\cdot10^{-15}$
\end{tabular}
\end{center}

Следовательно, уже примерно $2.52\cdot10^7$ проходов воспроизводят
рабочий второй момент с машинной точностью. Это согласуется с образом
очень большого числа витков. Но проект пока знает веса только как
вещественные результаты термического минимума. Он не доказывает, что они
являются точными рациональными частотами конечного цикла.

\section{Что говорит действие петли Тёплица}

Для собственных winding-чисел существующее действие имеет вид
\begin{equation}
 S_{\rm loop}=\sum_a k_a^2.
\end{equation}
При фиксированной полной намотке оно предпочитает распределять winding по
каналам, а не концентрировать его в одном канале. Именно поэтому класс
пятнадцать реализован пятнадцатью единичными winding и девяноста нулями,
а не одним winding $15$.

Для $10^{10}$ оборотов один канал стоил бы $10^{20}$. Равномерное
распределение по $105$ коэффициентным каналам уменьшает минимум примерно
в $105$ раз. Это косвенно поддерживает многопроходную картину, но
коэффициентные каналы $q_0\mathbb C^{105}$ ещё не отождествлены с
последовательными пространственными витками одной нити. Такое
отождествление нельзя вводить задним числом.

\begin{nogocriterion}
Высокоточное рациональное приближение не является выводом точных чисел
проходов. Также абстрактный цикл позиций не является вложенной нитью
конечной толщины: он не запрещает геометрическое самопересечение,
пересоединение или наложение нескольких витков на одну траекторию.
\end{nogocriterion}

\section{Вердикт и следующий тест}

Проект уже содержит точный абсолютный счётчик проходов $N$. Поэтому
гипотеза одной нити существенно усиливается: $C_4$, $Z_3$ и $Z_6$ можно
читать как конечные тени целого winding-числа, а не как полный счётчик.
Один абстрактный неветвящийся цикл с заданными целыми кратностями всегда
существует.

Но родитель ещё не выбирает единственное циклическое слово и не вкладывает
его в пространство как трубку конечного сечения. Следующий гейт должен
проверить рамированное геометрическое вложение: можно ли уложить слово
проходов без самопересечений, сохранив моменты, голономию и тип узла.

Нулевая гипотеза: существующие хопфова линия, порядок оператора $N$ и
тетраэдральный каркас канонически задают рамированное вложение одного
цикла.

Стоп-критерий: если порядок проходов или правило разнесения соседних
витков требует нового произвольного кода, толщины или матрицы
пересоединения, геометрическая неразрывность нити не выведена.

\begin{protocol}
\begin{enumerate}
 \item абсолютный целый счётчик проходов: \textbf{$N$};
 \item единичный шаг: \textbf{$U$};
 \item $[N,U^k]=kU^k$: \textbf{да};
 \item конечная голономия различает абсолютный $k$: \textbf{нет};
 \item один абстрактный цикл для заданных $K_a$: \textbf{существует};
 \item ветвления и концы в этом цикле: \textbf{отсутствуют};
 \item рабочие веса допускают точное конечное целочисленное чтение:
 \textbf{не доказано};
 \item приближение рабочего $Q$ до $3.13\cdot10^{-15}$:
 \textbf{$25185499$ проходов};
 \item единственное циклическое слово выведено: \textbf{нет};
 \item геометрическое вложение конечной толщины выведено: \textbf{нет};
 \item одна глобальная нить опровергнута: \textbf{нет};
 \item одна глобальная нить кинематически переоткрыта: \textbf{да};
 \item следующий тест: \textbf{рамированное winding-вложение}.
\end{enumerate}
\end{protocol}

Машинный сертификат сохранён в
\path{s2t_v6_single_thread_c4_suspension_parent_gate_results.json}.

\begin{thebibliography}{9}
\bibitem{single-thread-bourne-kellendonk-rennie2017}
C.~Bourne, J.~Kellendonk, A.~Rennie,
\emph{The $K$-theoretic bulk--edge correspondence for topological
insulators},
Annales Henri Poincar\'e 18 (2017), 1833--1866; arXiv:1604.02337.

\bibitem{single-thread-arici-mesland2020}
F.~Arici, B.~Mesland,
\emph{Toeplitz extensions in noncommutative topology and mathematical
physics},
arXiv:1911.05823.
\end{thebibliography}